\chapter{Durchführung} 
\label{ch:Durchfuehrung}

\section{Messverfahren}
\label{sec:Messverfahren}
Die Messung des radioaktiven Zerfalls erfolgt mit einem Geiger-Müller-Zählrohr. Dieses ist ein Messgerät für ionisierende Strahlung. Tritt ein ionisierendes Teilchen in das Zählrohr ein, so werden Gasteilchen ionisiert. Die entstehenden freien Elektronen werden durch die anliegende Spannung beschleunigt und können weitere Gasatome ionisieren. Dadurch steigt die Zahl an freien Elektronen schnell an, was als Elektronenlawine bezeichnet wird. Diese erzeugt einen Strompuls, der verstärkt wird und dann von einem Zähler detektiert werden kann. Damit die Elektronenlawine wieder abklingt, ist dem Gas ein Löschgas zugefügt.
Die Anzahl der ionisierten Gasatome hängt von der Spannung ab, mit der das Zählrohr betrieben wird. Da energiereichere einfallende Teilchen mehr Gasatome ionisieren, ermöglicht dieser Zusammenhang Rückschlüsse auf die Energie der einfallenden Teilchen basierend auf der Amplitude des Strompulses. Da in diesem Versuch jedoch nur die Anzahl an Zerfällen bestimmt werden soll, sollte das Zählverhalten der Apparatur unabhängig von der Energie der einfallenden Teilchen sein. Daher wird die Spannung so eingestellt, dass im Plateaubereich gemessen wird.
Zur Bestimmung der Plateauspannung $U_0$ werden die Countraten in Abhängigkeit von der Spannung $U$ am Zählrohr vermessen. Anschließend wird $U_0$ in der Mitte des Plateaubereichs gewählt, um sicherzustellen, dass der Zähler unabhängig von der Energie der einfallenden Teilchen zählt. Für die folgenden Messungen wird typischerweise eine Spannung von etwa $U_0 = 560\,\mathrm{V}$ gewählt.
Die eigentliche Messung erfolgt durch Zählen der Zerfälle in definierten Zeitintervallen. Dabei werden sowohl Messungen mit großer mittlerer Ereigniszahl als auch mit kleiner mittlerer Ereigniszahl durchgeführt, um den Übergang von der Poisson- zur Gauß-Verteilung zu untersuchen. Für die Messungen mit großer Ereigniszahl wird das radioaktive Präparat näher an das Zählrohr gebracht, während für die Messungen mit kleiner Ereigniszahl der Abstand vergrößert und die Torzeit reduziert wird.

\section{Versuchsaufbau}
\label{sec:Versucshaufbau}
Für die Durchführung des Versuchs werden folgende Materialien und Geräte benötigt: Ein Geiger-Müller-Zählrohr mit Betriebsgerät, ein externer Impulszähler, ein Personal Computer zur Datenerfassung und -auswertung, eine Präparathalterung mit Bleischirmung sowie ein radioaktives Präparat (typischerweise $^{60}\mathrm{Co}$).
Der Versuchsaufbau besteht aus dem Geiger-Müller-Zählrohr, welches über ein Kabel mit dem Betriebsgerät verbunden ist. Das Betriebsgerät liefert die notwendige Hochspannung für den Betrieb des Zählrohrs und wandelt die detektierten Ionisationsereignisse in elektrische Signale um. Ein externer Impulszähler erfasst die Anzahl der Detektionen in definierten Zeitintervallen. Der Personal Computer dient zur Steuerung des Experiments, zur Datenerfassung und zur anschließenden Auswertung der Messdaten mittels geeigneter Software.
Das radioaktive Präparat wird in einer speziellen Halterung positioniert, die eine präzise Einstellung des Abstands zum Zählrohr ermöglicht. Eine Bleischirmung sorgt dafür, dass nur die gewünschte Strahlung in Richtung des Zählrohrs gelangt und Streustrahlung minimiert wird. Die gesamte Anordnung wird auf einer stabilen Unterlage montiert, um Vibrationen und unbeabsichtigte Bewegungen während der Messungen zu vermeiden.

\subsubsection*{Skizze des Versuchsaufbaus}
\label{sec:Skizze}